% ============================================================
% Come si usano, argomento 9: funzioni di più variabili, covarianza e correlazione
% ============================================================
\argomento{Funzioni di più variabili, covarianza e correlazione}{Formule a p.~\pageref*{f:pmffunzione2}}[Covarianza
  e correlazione]

\begin{scelta}{Come scegliere la formula}
\renewcommand{\arraystretch}{1.25}
\begin{tabular}{@{}P{0.55\linewidth}P{0.39\linewidth}@{}}
  \toprule
  \textbf{Se il testo\ldots} & \textbf{si usa\ldots} \\
  \midrule
  chiede la pmf di $Z = g(X, Y)$ & la somma sulle coppie che danno $z$ \vv{pmffunzione2} \\
  chiede solo la media di $Z = g(X, Y)$ & \vv{valoremedio2}; se $g$ è lineare basta la linearità \vv{linearita} \\
  chiede la media di una somma di variabili, anche dipendenti & la linearità \vv{linearita} \\
  rende facile la media sapendo $Y$ & la media in due tempi \vv{mediacondizionata} \\
  chiede quanto due variabili «si muovono insieme» & covarianza \vv{covarianza} e, per un numero puro,
    $\rho$ \vv{rho} \\
  chiede la varianza di una somma o di una differenza & \vv{varsomma}: attenzione al termine con la covarianza \\
  \bottomrule
\end{tabular}

Per covarianza e $\rho$ conviene preparare, in quest'ordine: le marginali, $E[X]$, $E[Y]$, $E[X^2]$, $E[Y^2]$, le
varianze, $E[XY]$.
\end{scelta}

\scheda{pmffunzione2}{
  \usoquando{Si vuole la pmf di una variabile ottenuta combinando $X$ e $Y$: una somma, un prodotto, una formula
    qualsiasi.}
  \usocome{(1) Si calcola $g(x, y)$ in ogni casella della congiunta. (2) Per ogni valore $z$ si sommano le
    probabilità delle caselle che lo producono.}
  \usoesempio{Due bit con $p(0,0) = \frac{1}{3}$, $p(0,1) = \frac{2}{9}$, $p(1,0) = \frac{1}{9}$, $p(1,1) =
    \frac{1}{3}$ e $Z = 3X^2 + Y$: le quattro caselle danno $0, 1, 3, 4$, tutti diversi, quindi $p_Z = \left(\frac{1}{3},
    \frac{2}{9}, \frac{1}{9}, \frac{1}{3}\right)$. Con $X, Y \in \{-1, 1\}$, invece, $Z = 3 + Y$ e due caselle finiscono
    in ogni valore: le loro probabilità si sommano.}
}

\scheda{valoremedio2}{
  \usoquando{Serve la media di $g(X, Y)$ e non la sua pmf: per esempio $E[XY]$, $E[(X - Y)^2]$.}
  \usocome{Si pesa il valore di $g$ in ogni casella con la probabilità della casella e si somma su tutta la tabella.}
  \usoesempio{Con i due bit dell'esempio precedente, $E[3X^2 + Y] = 0 \cdot \frac{1}{3} + 1 \cdot \frac{2}{9} + 3 \cdot
    \frac{1}{9} + 4 \cdot \frac{1}{3} = \frac{17}{9}$.}
}

\scheda{linearita}{
  \usoquando{Si vuole la media di una somma (o di una combinazione lineare) di variabili: spesso conviene scrivere la
    variabile che interessa come somma di pezzi semplici.}
  \usoattenzione{La linearità \emph{non} richiede l'indipendenza: vale sempre. Per la varianza di una somma, invece,
    la dipendenza conta (\vv{varsomma}).}
  \usoesempio{Le teste in $N$ lanci sono la somma di $N$ variabili di Bernoulli, ognuna con media $p$:
    $E[X_N] = p + p + \cdots + p = N p$, in una riga.}
}

\scheda{mediacondizionata}{
  \usoquando{La media di $g(X, Y)$ è facile da calcolare quando il valore di $Y$ è noto: prima la si calcola per ogni
    $y$, poi si fa la media su $Y$.}
  \usocome{(1) Per ogni $y$ si trova la condizionata $p_{X \mid Y}(x \mid y)$ e la media $h(y)$. (2) Si pesano le
    $h(y)$ con $p_Y(y)$. I ruoli di $X$ e $Y$ si possono scambiare.}
  \usoesempio{Con i due bit di prima: $E[3X^2 + Y \mid Y = 0] = 3 \cdot \frac{1}{4} = \frac{3}{4}$ ed
    $E[3X^2 + Y \mid Y = 1] = 3 \cdot \frac{3}{5} + 1 = \frac{14}{5}$; quindi
    $E[Z] = \frac{3}{4} \cdot \frac{4}{9} + \frac{14}{5} \cdot \frac{5}{9} = \frac{17}{9}$, come prima.}
}

\scheda{correlazione}{
  \usoquando{Come passo intermedio per la covarianza; e quando una delle due variabili ha media nulla, perché in
    quel caso correlazione e covarianza coincidono.}
  \usocome{Si sommano i prodotti $x\,y$ pesati con la congiunta: le caselle con $x = 0$ o $y = 0$ non contano.}
  \usoesempio{Due bit: $E[XY] = 1 \cdot 1 \cdot p(1, 1) = \frac{1}{3}$.}
}

\scheda{covarianza}{
  \usoquando{Si vuole sapere se $X$ e $Y$ tendono a crescere insieme (covarianza positiva) o in direzioni opposte
    (negativa); serve per la varianza di una somma e per $\rho$.}
  \usocome{Quasi sempre con $E[XY] - E[X]\, E[Y]$.}
  \usoattenzione{Dipende dalle unità di misura: per confrontare coppie diverse si usa $\rho$ (\vv{rho}).}
  \usoesempio{Due bit con $E[X] = \frac{4}{9}$, $E[Y] = \frac{5}{9}$, $E[XY] = \frac{1}{3}$:
    $COV(X, Y) = \frac{1}{3} - \frac{20}{81} = \frac{7}{81}$.}
}

\scheda{incorrelate}{
  \usoquando{Per semplificare: con variabili incorrelate (in particolare indipendenti) $E[XY]$ si spezza nel prodotto
    delle medie e le varianze si sommano.}
  \usoattenzione{Incorrelate non vuol dire indipendenti. $X$ uniforme su $\{-1, 0, 1\}$ e $Y = X^2$: $Y$ è determinata
    da $X$, eppure $E[XY] = E[X^3] = 0 = E[X]\, E[Y]$.}
  \usoesempio{Una covarianza diversa da zero basta a dire che le variabili sono dipendenti: i due bit di prima, con
    $COV = \frac{7}{81}$, lo sono.}
}

\scheda{cauchy}{
  \usoquando{Come controllo dei conti: se una covarianza supera in valore assoluto $\sigma_X \sigma_Y$, c'è un errore.
    È anche il motivo per cui $\rho$ sta fra $-1$ e $1$.}
  \usoesempio{Due bit con $\sigma_X^2 = \sigma_Y^2 = \frac{20}{81}$: la covarianza non può superare $\frac{20}{81}$, e
    infatti vale $\frac{7}{81}$.}
}

\scheda{rho}{
  \usoquando{Si vuole un numero puro, senza unità di misura, che dica quanto due variabili sono legate
    linearmente, o si devono confrontare coppie diverse.}
  \usocome{Si divide la covarianza per il prodotto delle deviazioni standard. Il segno dice la direzione, il valore
    assoluto la forza.}
  \usoattenzione{$\rho$ vicino a 0 non esclude una dipendenza non lineare (come $Y = X^2$).}
  \usoesempio{Due bit: $\rho = \frac{7/81}{20/81} = 0{,}35$. Con $X, Y \in \{-1, 1\}$ e $p(-1,-1) = \frac{1}{4}$,
    $p(-1,1) = \frac{1}{2}$, $p(1,\pm 1) = \frac{1}{8}$: $\rho = -\frac{1}{\sqrt{45}} \simeq -0{,}149$, più debole e di
    segno opposto.}
}

\scheda{varsomma}{
  \usoquando{Si vuole la varianza di una somma, di una differenza o di una combinazione lineare di variabili.}
  \usocome{Si sommano le varianze con i quadrati dei coefficienti e si aggiunge il termine con la covarianza; con
    variabili indipendenti quel termine è nullo.}
  \usoattenzione{La varianza di una differenza di variabili indipendenti è la \emph{somma} delle varianze, non la
    differenza: $\sigma_{X-Y}^2 = \sigma_X^2 + \sigma_Y^2$.}
  \usoesempio{La somma di due dadi indipendenti ha varianza $\frac{35}{12} + \frac{35}{12} = \frac{35}{6} \simeq
    5{,}83$. La binomiale è la somma di $N$ Bernoulli indipendenti: $\sigma^2 = N p\,(1 - p)$.}
}
